\documentclass[12pt]{article}
\usepackage{geometry}
\geometry{a4paper, margin=1in}
\usepackage{fancyhdr}
\usepackage{titlesec}
\usepackage{hyperref}
\usepackage{graphicx}
\usepackage{amsmath}
\usepackage{listings}

% Title Page
\title{
    \vspace{0.2in}
    \Huge \textbf{Maulana Abul Kalam Azad University of Technology} \\
    \vspace{0.5in} % Reduced the vertical space
    \begin{center}
        \includegraphics[width=0.6\textwidth]{Pic.png} % Add your logo image here
    \end{center}
    \vspace{0.5in}
    \Huge\textbf{\textit{Lab Notebook}} \\
    \vspace{0.5in}
    \newpage
    \Large \textbf{\textit{GROUP MEMBERS}} \\
    \vspace{0.5in}
}
\date{}

\begin{document}
\maketitle

% Group Members Information Table
\begin{center}
    \begin{tabular}{|c| c | c | c |}
    \hline
    \textbf{SL NO}&\
    \textbf{Name} & \textbf{Roll Number} & \textbf{Department} \\
    \hline
    1 & TAMAGHNA GHOSH &30001223025 & BCA  \\
    \hline
    2& SRITI MAJUMDER & 30001223078 & BCA \\
    \hline
    3& SOUMYAJIT KHAN & 30001223019 & BCA   \\
    \hline
    4 & HIMADRI BERA & 30001223009 & BCA \\
    \hline
    5 & MOHIMA ISLAM & 30001223049 & BCA \\
    \hline
    \end{tabular}
\end{center}
\newpage
\begin{center}
    \begin{tabular}{|c|  c | c|c|c| }
    \hline
    \textbf{SL NO}&\
    \textbf{INDEX} & \textbf{DATE}&\textbf{ENTRY BY} &\textbf{PAGE NO} \\
    \hline
    1 & My First \LaTeX & 30/8/24&TAMAGHNA GHOSH & 4-7\\
    \hline
    2&Git Branching and Merging&1/9/24& TAMAGHNA GHOSH &8 \\
    \hline
    3&  MY CV & 3/9/24 & --&\\
    \hline
    4&   -- &-- &--&--\\
    5& GIT ASSIGNMENT STT & 16/08/24 & HIMADRI BERA &--\\
    6&--&--&--&--\\
    \end{tabular}
\end{center}

\newpage

% Lab Notebook Entries
\section{Lab Notebook Entry by Member 1}
\subsection*{Date: 08-09-2024}

\begin{flushright}
\textbf{Name:} TAMAGHNA GHOSH \\
\textbf{Semester:} 2nd \\
\textbf{Year:} I \\
\end{flushright}

\begin{center}
\Huge \textbf{ASSIGNMENT: GIT}
\end{center}

I am really grateful to my professor for assigning such a wonderful project, which has helped us gain valuable work experience. I also appreciate my peers for choosing me as the group leader, and I hope to meet their expectations in this role.


\newpage

\section{MY FIRST LATEX DOCUMENT}
\subsection*{Date: 30-08-2024}

\begin{center}
    \includegraphics[width=0.2\textwidth]{tamaghna.jpg}\\[1cm]
    
\end{center}

\section*{Introduction}
Hello !! My name is Tamaghna Ghosh. I am currently a 1st year undergrad at MAKAUT University doing my BCA.

\subsection*{Interests \& Hobbies}
My hobbies are drawing and a bit of singing. I love to paint and have a nigh interest towards graphic design. Some of my other hobies are :

\begin{itemize}
    \item Playing batminton
    \item hanging out with friends
        \begin{itemize}
            \item Usually catching up after classes
        \end{itemize}
    \item Doing Coding
\end{itemize}

\subsection*{Favorite Quotations}
\begin{enumerate}
    \item I think I can. I know I can. – Jennifer Wittwer
    \item Only I can change my life. No one can do it for me. – Carol Burnett
\end{enumerate}


\section*{2 Mathematics}

\subsection*{2.1 Mathematics and Me}
Mathematics has always been a very favourite subject of mine. It has always interested me and motivated me to believe in the fact that every problem is solveable.

\subsection*{2.2 Mathematical Notation}
I choose the number 2004 

\begin{enumerate}
    \item Superscripts, subscripts, and Greek letters
    \begin{itemize}
        \item \(2004\)
        \item \(\mu^{2004}\)
        \item \(\pi_{2004}\)
        \item \(2004^{\alpha}\)
        \item \(2004_m\)
        \item \(\tau_{2004}\)
        \item \(\log_{2004}\)
        \item \(\ln 2004\)
        \item \(0 < x \leq 2004\)
        \item \(\sum_{i=1}^{2004} x_i\)
    \end{itemize}
    \item Roots, fractions, and displaystyle
    \begin{itemize}
        \item \(\sqrt{2004}\)
        \item \(\sqrt[3]{2004}\)
        \item \(\frac{2004}{2}\)
        \item \(\frac{\sqrt{2004}}{2}\)
        \item \(\frac{1}{\sqrt{2004}}\)
    \end{itemize}
    \item Delimiters
    \begin{itemize}
        \item \(\left( \frac{2004}{2} \right)\)
        \item \(\left[ \frac{2004}{2} \right]\)
        \item \(\left\{ \frac{2004}{2} \right\}\)
    \end{itemize}
\end{enumerate}

\section*{Tables and Continuous Functions}
\begin{enumerate}
    \item[(a)] \(\lim_{x \to 3} \frac{x^2 - 9}{x - 3} = ?\)
    \item[(b)] \(\lim_{x \to -2} \frac{x^2 + x - 6}{x + 2} = ?\)
    \item[(c)] \(\lim_{h \to 0} \frac{(3+h)^2 - 9}{h} = ?\)
\end{enumerate}

\section*{Functions \& Formulas}
\begin{enumerate}
    \item[(a)] \(f(x) = x^3 - x\)
    \begin{enumerate}
        \item[i.] \(f'(x) = ?\)
        \item[ii.] \(f''(x) = ?\)
    \end{enumerate}
    \item[(b)] \(g(x) = \frac{1}{1+x^2}\)
    \begin{enumerate}
        \item[i.] \(g'(x) = ?\)
    \end{enumerate}
\end{enumerate}

\section*{Fundamental Theorem of Calculus, Part I}
If \(F\) is an antiderivative for \(f\) on an interval \(I\), then for every continuous function on \(I\),
\[ \int_a^b f(x)\, dx = F(b) - F(a). \]

\section*{Fundamental Theorem of Calculus, Part II}
If \(f\) is continuous on \([a,b]\), then the function defined by
\[ F(x) = \int_a^x f(t)\, dt \]
for all \(x\) in \([a,b]\) is continuous on \([a,b]\) and differentiable on \((a,b)\), and \(F'(x) = f(x)\).

\section*{Example}
Method vs use of F.T.C., part 2 in the step-by-step solution.
\[ \int_0^{1972} \frac{dx}{1+x^2} = [\arctan x]_0^{1972} = \ldots \]

\section*{Taylor Series}
Convergent geometric series since \(|r| < 1\).
\[ e^{-t/12} \sin t = \sum_{n=0}^\infty (-1)^n \frac{t^{n+1}}{(12)^n (n+1)!} \]
\newpage
\section{Git Branching and Merging}
\subsection*{Date: 08-09-2024}
\title{Git Branching and Merging{} Document}
Assignment: Git Branching and Merging

Objective: Demonstrate proficiency in Git branching, merging, and conflict resolution.

Task step by step:

1. Create a new repository on GitHub called "git-advanced".
2. Clone the repository to your local machine.
3. Create a new branch called "feature-1" and switch to it.
4. Create a new file called "shared.txt" and add the following text:

This is a shared file.
Line 1: Original text.
Line 2: Original text.

5. Stage and commit the file with a meaningful message.
6. Push the branch to GitHub.
7. Create another branch called "feature-2" and switch to it.
8. Checkout the "shared.txt" file from the main branch (to ensure you're starting with the same file).
9. Modify the "shared.txt" file by changing the second line to:

Line 2: Modified text in feature-2.

10. Stage and commit the changes with a meaningful message.
11. Push the branch to GitHub.
12. Switch back to the "feature-1" branch.
13. Modify the "shared.txt" file by changing the second line to:

Line 2: Modified text in feature-1.

14. Stage and commit the changes with a meaningful message.
15. Push the branch to GitHub.
16. Merge "feature-1" into the main branch.
17. Merge "feature-2" into the main branch, which should introduce a conflict.
18. Resolve the conflict by editing the file and committing the resolution.
19. Push the updated main branch to GitHub.
20. Delete the "feature-1" and "feature-2" branches.

Deliverables:
You need to upload only ONE pdf file consists of the below things:
A screenshot of your GitHub repository showing the commit history and branching.
A screenshot of your local machine showing the Git log and conflict resolution.
A brief write-up of your experience with Git branching and merging.
Grading criteria:
Correct creation and management of branches.
Accurate merging and conflict resolution.
Clear and meaningful commit messages.
Proper deletion of branches.
\newpage


% Lab Notebook Entries
\section{Lab Notebook Entry by Member 2}
\subsection*{Date: 08-09-2024}

\begin{flushright}
\textbf{Name:} SRITI MAJUMDER \\
\textbf{Semester:} 2nd \\
\textbf{Year:} I \\
\end{flushright}

\begin{center}
\Huge \textbf{ASSIGNMENT: GIT}
\end{center}

I am immensely grateful to my teacher, Ayan Sir, for assigning this remarkable group project, which has provided us with invaluable hands-on experience. I also want to extend my appreciation for the unwavering support and teamwork throughout this endeavor. Additionally, I have able to learn how to  resolved minor issues in my  code to enhance its functionality.

\newpage
\section{MY FIRST LATEX DOCUMENT}
\subsection*{Date: 30-08-2024}
\newpage
% Lab Notebook Entries
\section{Lab Notebook Entry by Member 3}
\subsection*{Date: 08-09-2024}

\begin{flushright}
\textbf{Name:} SOUMYAJIT KHAN \\
\textbf{Semester:} 2nd \\
\textbf{Year:} I \\
\end{flushright}

\begin{center}
\Huge \textbf{ASSIGNMENT: GIT}
\end{center}

I am deeply thankful to my teacher, Ayan Sir, for giving us this exceptional group project. It has been an incredible opportunity for hands-on learning and skill development. I also want to express my gratitude for the consistent support and collaborative effort throughout this journey. Furthermore, I've gained valuable experience in troubleshooting and refining my code to improve its functionality.


\newpage
% Lab Notebook Entries
\section{GIT ASSIGNMENT STT}
\subsection*{Date: 08-09-2024}

 \begin{lstlisting}[language=JAVA, caption={STT}]
import java.awt.*;
import java.awt.event.*;
import java.util.Random;
public class SymbolApp extends Frame implements ActionListener {
    private Label[] symbolLabels = new Label[99];
    private Button submitButton;
    private String specialSymbol;
    private String selectedSymbol;
    public SymbolApp() {
        // Generate a random special symbol
        Random rand = new Random();
        specialSymbol = Character.toString((char) (rand.nextInt(94) + 33)); // Random ASCII character from 33 to 126
        selectedSymbol = specialSymbol;
        // Setting up the main frame
        setLayout(new BorderLayout());
        setSize(800, 700);
        setTitle("Symbol App");
        // Adding instruction message
        TextArea instruction = new TextArea(
            "Think of any two digit number. Now reverse it and find the difference of them.\n" +
            "Now find the number you got and remember the symbol from the panel below.\n" +
            "Don't tell me, I'll read your mind! Hit the below button when you are ready to see the magic!",
            5, 60, TextArea.SCROLLBARS_NONE);
        instruction.setEditable(false);
        instruction.setFont(new Font("Arial", Font.PLAIN, 16));
        add(instruction, BorderLayout.NORTH);
        // Panel for symbols
        Panel symbolPanel = new Panel(new GridLayout(11, 9));
        for (int i = 0; i < 99; i++) {
            String symbol = (i % 9 == 0) ? specialSymbol : Character.toString((char) (33 + (i % 94)));
            symbolLabels[i] = new Label(i + ": " + symbol); // Numbering symbols
            symbolLabels[i].setAlignment(Label.CENTER);
            symbolPanel.add(symbolLabels[i]);
        }
        add(symbolPanel, BorderLayout.CENTER);
        // Panel for submit button
        Panel controlPanel = new Panel(new FlowLayout());
        submitButton = new Button("Chin Tapak Tapak Dum Dum");
        submitButton.addActionListener(this);
        controlPanel.add(submitButton);
        add(controlPanel, BorderLayout.SOUTH);
        // Setting up the window close event
        addWindowListener(new WindowAdapter() {
            public void windowClosing(WindowEvent we) {
                System.exit(0);
            }
        });
        setVisible(true);
    }
    public void actionPerformed(ActionEvent ae) {
        // Clear the current content and display the selected symbol
        if (ae.getSource() == submitButton) {
            removeAll();
            setLayout(new BorderLayout());
            Label resultLabel = new Label(selectedSymbol, Label.CENTER);
            resultLabel.setFont(new Font("Arial", Font.BOLD, 50));
            add(resultLabel, BorderLayout.CENTER);
            validate();
            repaint();
        }
    }
    public static void main(String[] args) {
        new SymbolApp();
    }
} 
\end{lstlisting}
\begin{flushright}
\textbf{Name:} HIMADRI BERA \\
\textbf{Semester:} 2nd \\
\textbf{Year:} I \\
\end{flushright}

\begin{center}
\Huge \textbf{ASSIGNMENT: GIT}
\end{center}

I'm incredibly thankful to my professor for assigning such a fantastic project, which has given us invaluable hands-on experience. I'm also grateful to my teammates for their collaboration, and I hope to contribute meaningfully to the team's success.

\newpage
\section{Lab Notebook Entry by Member 4}
\subsection*{Date: 08-09-2024}
\newpage


% Lab Notebook Entries
\section{Lab Notebook Entry by Member 5}
\subsection*{Date: 08-09-2024}

\begin{flushright}
\textbf{Name:} MOHIMA ISLAM \\
\textbf{Semester:} 2nd \\
\textbf{Year:} I \\
\end{flushright}

\begin{center}
\Huge \textbf{ASSIGNMENT: GIT}
\end{center}

I express my gratitude towards my teacher for giving the opportunity to do this assignment. This project has enhanced my experience, skill and also gave me a chance to do teamwork. I am also thankful to my teammates for their support and I'm praying for the best result. Thank you.
\newpage
\section{Lab Notebook Entry by Member 5}
\subsection*{Date: 08-09-2024}


\end{document}
